\section{Investigación sobre infraestructura óptica y servicios}
\label{sec:investigacion}

\subsection{REFEFO y ARSAT}
\label{sec:refefo}
% Consigna 1.1:
% - Significado de la sigla y objetivos de la red.
% - Desarrollo, cobertura y servicios que ofrece.
% - Función de la REFEFO como infraestructura de transporte y conectividad mayorista.
% - Relación con los operadores que brindan servicios de acceso a los usuarios.
% - Actividades de ARSAT vinculadas con fibra óptica, televisión digital y servicios satelitales.
% - Indicar fuentes y FECHA DE CONSULTA.
% Fuente útil: material de cátedra, "Libro Cap 1 a 5 parte II", sección 1.7.
\pendiente{síntesis sobre la REFEFO y ARSAT, con fuentes y fecha de consulta.}

\subsection{Fundamentos de fibra óptica}
\label{sec:fundamentos}
% Consigna 1.2: síntesis integrada de los 5 videos indicados por la cátedra y
% explicación de qué conceptos son relevantes para el diseño FTTH y la selección de componentes.
\pendiente{síntesis integrada de los videos y su relevancia para el diseño FTTH.}

\subsection{Cables submarinos}
\label{sec:submarinos}
% Consigna 1.3: mapa de cables submarinos con llegada a Las Toninas.
% - Sistemas representados, destinos principales y estado informado.
% - Función en la conectividad internacional de Argentina y relación con las redes terrestres.
% Fuente útil: material de cátedra, "Libro Cap 1 a 5 parte II", sección 1.8.
\pendiente{sistemas que llegan a Las Toninas, destinos, estado y su rol en la conectividad.}
